\begin{notechapter}
\chapter{Object Detection and Instance Segmentation}\label{ch:detection-segmentation}
\chaptersubtitle{Boxes, proposals, RoI features and masks: the R-CNN family and its successors}

\begin{chapterintro}
Classification says \emph{what} is in an image. Detection also says \emph{where}, with a box around
each object; instance segmentation gives each object's exact pixels. These tasks need
new ingredients on top of a CNN: a way to propose or enumerate candidate regions, to
extract features for a region of any size, to score and refine boxes, and to remove
duplicates. This chapter builds them up through the R-CNN family, which culminates in
Mask R-CNN, and then describes one-stage alternatives. You will learn
\begin{itemize}[nosep]
  \item the vision tasks from classification to panoptic segmentation;
  \item intersection over union, non-maximum suppression and the box-regression
        parametrization, with worked numbers;
  \item R-CNN, Fast R-CNN and Faster R-CNN: proposals, RoI pooling, anchors and the region
        proposal network;
  \item RoI Align and bilinear interpolation, and the mask branch of Mask R-CNN;
  \item one-stage detectors (YOLO, SSD, RetinaNet), YOLACT and SOLO, and how detectors are
        evaluated.
\end{itemize}
\prerequisites{\cref{ch:convolutional-layers,ch:cnn-architectures,ch:fully-convolutional};
cross-entropy and the focal loss (\cref{surrogate-losses:sec:beyond}).}
\end{chapterintro}

\section{The tasks}\label{detection-segmentation:sec:tasks}

\begin{itemize}
  \item \textbf{Classification}: one label per image.
  \item \textbf{Object detection}: for each object, a class and an axis-aligned \emph{bounding box}
        $(x_1, y_1, x_2, y_2)$ (corners) or $(x, y, w, h)$ (centre and size), with a confidence score.
  \item \textbf{Semantic segmentation}: a class per pixel, without separating objects: two adjacent
        cars form one ``car'' region (\cref{fully-convolutional:sec:segmentation}).
  \item \textbf{Instance segmentation}: a class and a pixel mask for each object: ``car 1'' and ``car 2'' have
        separate masks, which may be adjacent.
  \item \textbf{Panoptic segmentation} (Kirillov et al., CVPR, 2019): every pixel gets a class, and pixels of
        countable objects (\emph{things}: cars, people) also get an instance identifier; amorphous
        regions (\emph{stuff}: sky, road) do not.
\end{itemize}
The standard benchmark is COCO (Lin et al., \emph{Microsoft COCO: Common Objects in
Context}, ECCV, 2014), with $80$ object classes and instance masks.

\section{Geometry of boxes}\label{detection-segmentation:sec:boxes}

\begin{definition}[Intersection over union]\label{detection-segmentation:def:iou}
For two boxes (or masks) $A$ and $B$, $\mathrm{IoU}(A,B) = \abs{A\cap B}/\abs{A\cup B}\in[0,1]$, with $\abs{\cdot}$ the area.
\end{definition}

For boxes, the intersection is the box $[\max(x_1, x_1'), \min(x_2, x_2')]\times[\max(y_1, y_1'), \min(y_2, y_2')]$ when both
intervals are non-empty, and $\abs{A\cup B} = \abs A + \abs B - \abs{A\cap B}$.

\begin{example}[IoU]\label{detection-segmentation:ex:iou}
$A = (0,0,4,4)$ and $B = (2,1,6,5)$ in corner form: the intersection is $[2,4]\times[1,4]$, of area $6$; the
union has area $16 + 16 - 6 = 26$; $\mathrm{IoU} = 6/26 = 0.23$ (\cref{detection-segmentation:fig:iou}). Two boxes of the
same size shifted by a quarter of their width already have $\mathrm{IoU} = 0.6$: the measure is strict.
\end{example}

\begin{figure}[ht]
\centering
\begin{tikzpicture}[scale=0.55]
  \draw[step=1, gray!20] (-0.5,-0.5) grid (6.5,5.5);
  \fill[bookexample!30] (2,1) rectangle (4,4);
  \draw[very thick, bookaccent] (0,0) rectangle (4,4) node[pos=0, below left, font=\small] {$A$};
  \draw[very thick, bookpitfall] (2,1) rectangle (6,5) node[pos=1, above right, font=\small] {$B$};
  \node[font=\small] at (3,2.5) {$6$};
  \node[font=\small, align=left, anchor=west] at (7,2.5) {$\abs A = \abs B = 16$\\$\abs{A\cap B} = 6$\\$\mathrm{IoU} = \frac6{26} = 0.23$};
\end{tikzpicture}
\caption{Intersection over union of two boxes.}
\label{detection-segmentation:fig:iou}
\end{figure}

\paragraph{Non-maximum suppression.} A detector typically fires several overlapping boxes
on the same object. NMS keeps the best one: sort the boxes of a class by score; repeatedly keep
the highest-scoring remaining box and discard the remaining boxes whose IoU with it exceeds
a threshold $\tau$ (often $0.5$).

\begin{example}[NMS]\label{detection-segmentation:ex:nms}
Boxes $b_1 = (0,0,4,4)$, $b_2 = (1,0,5,4)$, $b_3 = (6,6,9,9)$, $b_4 = (0,1,4,5)$ with scores $0.9, 0.8, 0.7, 0.6$, $\tau = 0.5$.
Keep $b_1$; $\mathrm{IoU}(b_1, b_2) = 12/20 = 0.6$ and $\mathrm{IoU}(b_1, b_4) = 0.6$, so $b_2$ and $b_4$ are removed; $b_3$ does not overlap and is
kept. Result: $\{b_1, b_3\}$. NMS can wrongly remove a true object heavily overlapped by another one
(a crowd); Soft-NMS lowers scores instead of removing boxes.
\end{example}

\paragraph{Box regression.} A detector refines a reference box $(x_a, y_a, w_a, h_a)$ (a proposal or an
anchor) into a predicted box by predicting four numbers $t = (t_x, t_y, t_w, t_h)$, decoded as
\[
  x = x_a + w_at_x,\qquad y = y_a + h_at_y,\qquad w = w_ae^{t_w},\qquad h = h_ae^{t_h} .
\]
Offsets are relative to the reference size, so that small and large boxes are treated alike,
and sizes are on a log scale, so that the predicted size is always positive (Girshick et al.,
2014). The training targets are obtained by inverting these formulas for the ground-truth
box, and the loss is a smooth $L^1$ (Huber) loss on $t$. For the anchor $(50, 50, 20, 40)$ and the
ground truth $(54, 46, 30, 40)$, the targets are $t = (0.2, -0.1, \log1.5, 0) = (0.2, -0.1, 0.405, 0)$.

\section{The R-CNN family}\label{detection-segmentation:sec:rcnn}

\paragraph{R-CNN} (Girshick, Donahue, Darrell and Malik, \emph{Rich Feature Hierarchies for
Accurate Object Detection and Semantic Segmentation}, CVPR, 2014) treats detection as
classification of regions:
\begin{enumerate}[nosep]
  \item about $2000$ region proposals per image from Selective Search (Uijlings et al., IJCV, 2013),
        a bottom-up grouping of superpixels;
  \item each proposal warped to $227\times227$ and passed through a CNN pretrained on ImageNet and
        fine-tuned (\cref{ch:transfer-learning}), giving a $4096$-dimensional feature vector;
  \item one linear SVM per class on these features (\cref{ch:svm}), plus box regression;
  \item NMS per class.
\end{enumerate}
It was a large jump in accuracy, but slow: $2000$ CNN passes per image.

\paragraph{Fast R-CNN} (Girshick, ICCV, 2015) runs the CNN once on the whole image and
extracts each proposal's features from the shared feature map by \emph{RoI pooling}: the
proposal is projected onto the feature map (divided by the stride, e.g.\ $16$), rounded to
integer cells, divided into a fixed grid of $7\times7$ bins (also rounded), and each bin is max-pooled.
Every proposal thus yields a $C\times7\times7$ tensor, whatever its size, which feeds fully connected
layers with two outputs: class scores (softmax over $K$ classes plus background) and
class-specific box regressions. Training is end to end with the multi-task loss
$\ell = \ell_{\mathrm{cls}} + \lambda\,\bbone[\text{not background}]\,\ell_{\mathrm{box}}$.

\paragraph{Faster R-CNN} (Ren, He, Girshick and Sun, \emph{Faster R-CNN: Towards Real-Time
Object Detection with Region Proposal Networks}, NeurIPS, 2015) replaces Selective
Search by a \emph{region proposal network} (RPN) sharing the backbone. At every position of the
feature map, $A$ \emph{anchors} (reference boxes of several scales and aspect ratios, e.g.\ $3\times3 = 9$) are
placed; a small fully convolutional head predicts for each anchor an \emph{objectness} score (object
or background, regardless of class) and a box regression $t$. Anchors with IoU above $0.7$ with a
ground-truth box are positive, below $0.3$ negative. The top proposals after NMS go to the
Fast R-CNN head. The whole detector is one network, two stages
(\cref{detection-segmentation:fig:faster}).

\begin{figure}[ht]
\centering
\begin{tikzpicture}[>=Stealth, op/.style={draw, rounded corners=2pt, fill=#1!15, font=\small, minimum height=7mm, align=center}]
  \node[op=gray] (img) at (0,0) {image};
  \node[op=bookaccent] (bb) at (2.4,0) {backbone\\(CNN, FPN)};
  \node[op=bookdef] (rpn) at (5.5,1.3) {RPN: objectness\\and boxes per anchor};
  \node[op=bookpitfall] (roi) at (5.5,-1.0) {RoI Align\\$C\times7\times7$ per proposal};
  \node[op=bookexample] (head) at (9.6,-0.2) {class scores\\box refinement};
  \node[op=bookexample] (mask) at (9.6,-1.8) {mask $28\times28$\\(Mask R-CNN)};
  \draw[->] (img) -- (bb);
  \draw[->] (bb.east) -- (rpn.west);
  \draw[->] (bb.east) -- (roi.west);
  \draw[->] (rpn.south) -- (roi.north) node[midway, right, font=\scriptsize] {proposals};
  \draw[->] (roi.east) -- (head.west);
  \draw[->] (roi.east) -- (mask.west);
\end{tikzpicture}
\caption{Faster R-CNN, and Mask R-CNN with its extra branch. The backbone runs once per
image; the heads run once per proposal on fixed-size RoI features.}
\label{detection-segmentation:fig:faster}
\end{figure}

\paragraph{Feature pyramids.} Objects come in all sizes. A feature pyramid network (Lin et
al., CVPR, 2017) adds a top-down path with lateral connections to the backbone, so that
every scale has semantically strong features; small objects are handled on fine levels,
large ones on coarse levels.

\section{Mask R-CNN and RoI Align}\label{detection-segmentation:sec:mask}

\paragraph{The quantization problem.} RoI pooling rounds the proposal and the bins to
integer cells of the feature map. At stride $16$, an error of half a cell is $8$ pixels in the image:
harmless for classification, damaging for pixel-accurate masks.

\paragraph{RoI Align} (He, Gkioxari, Dollár and Girshick, \emph{Mask R-CNN}, ICCV, 2017)
removes all rounding. The proposal is projected with real coordinates, divided into $7\times7$ bins
of real size, and in each bin a few regularly spaced sample points (e.g.\ $2\times2$) are evaluated by
\emph{bilinear interpolation} of the four surrounding feature cells, then averaged (or max-pooled).
At a point $(x,y) = (x_0 + \alpha, y_0 + \beta)$ with $\alpha,\beta\in[0,1)$ between cells,
\[
  f(x,y) = (1 - \alpha)(1 - \beta)f_{00} + \alpha(1 - \beta)f_{10} + (1 - \alpha)\beta f_{01} + \alpha\beta f_{11},
\]
where $f_{ab}$ is the feature at $(x_0 + a, y_0 + b)$. For $f_{00} = 1$, $f_{10} = 3$, $f_{01} = 5$, $f_{11} = 7$ and $(\alpha,\beta) = (0.25, 0.5)$:
$0.375 + 0.375 + 1.875 + 0.875 = 3.5$. Bilinear interpolation is differentiable in the features, so gradients flow
back to the backbone.

\paragraph{The mask branch.} Mask R-CNN adds to Faster R-CNN a small fully convolutional
head on the RoI features (often $14\times14$), which outputs $K$ masks of $28\times28$, one per class. Only the mask of
the ground-truth class is trained, with a per-pixel binary cross-entropy (sigmoid, not
softmax): masks of different classes do not compete, and classification is left to the box
head. At inference, the mask of the predicted class is resized to the box and thresholded at
$0.5$. The loss is $\ell = \ell_{\mathrm{cls}} + \ell_{\mathrm{box}} + \ell_{\mathrm{mask}}$. RoI Align alone improved mask accuracy by several
points of average precision.

\section{One-stage detectors and instance segmentation}\label{detection-segmentation:sec:onestage}

Two-stage detectors are accurate but run a head per proposal. One-stage detectors predict
classes and boxes densely, directly from the feature maps:
\begin{itemize}
  \item \textbf{YOLO} (Redmon, Divvala, Girshick and Farhadi, \emph{You Only Look Once},
        CVPR, 2016) divides the image into an $S\times S$ grid; each cell predicts a few boxes with
        confidences and class probabilities, in one forward pass, in real time.
  \item \textbf{SSD} (Liu et al., ECCV, 2016) predicts anchor-based boxes from several feature maps of
        different resolutions.
  \item \textbf{RetinaNet} (Lin, Goyal, Girshick, He and Dollár, ICCV, 2017) identified the main
        problem of dense detectors: about $100\,000$ anchors per image, almost all easy background,
        whose many small losses dominate the cross-entropy. The \emph{focal loss}
        (\cref{surrogate-losses:sec:beyond}) down-weights well-classified examples, and brought
        one-stage accuracy to the level of two-stage detectors.
\end{itemize}
Instance segmentation also has one-stage methods. \textbf{YOLACT} (Bolya, Zhou, Xiao and Lee,
ICCV, 2019) predicts a set of $k$ prototype masks for the whole image and, for each detection, $k$
coefficients; the instance mask is the sigmoid of the linear combination of the prototypes,
cropped to the box. \textbf{SOLO} (Wang, Kong, Shen, Jiang and Li, \emph{SOLO: Segmenting
Objects by Locations}, ECCV, 2020) avoids boxes: the image is divided into an $S\times S$ grid, and the
grid cell containing an object's centre predicts its class and its full-image mask.

\section{Evaluation}\label{detection-segmentation:sec:evaluation}

A predicted box is a true positive if its IoU with a not-yet-matched ground-truth box of the
same class is at least a threshold (predictions are matched greedily in order of decreasing
score); otherwise it is a false positive. Unmatched ground-truth boxes are false negatives.
Sweeping the score threshold traces a precision--recall curve (\cref{ch:measuring-performance});
its area is the \emph{average precision} (AP) of the class, and mAP averages over classes. PASCAL VOC
uses IoU $0.5$ (AP$_{50}$); COCO averages AP over IoU thresholds $0.5, 0.55,\dots,0.95$, which rewards precise
localization. Masks are evaluated in the same way with mask IoU.

\begin{chaptersummary}
\begin{itemize}[nosep]
  \item Detection outputs classes and boxes; semantic segmentation a class per pixel; instance
        segmentation a mask per object; panoptic segmentation both.
  \item IoU measures overlap; NMS keeps the best of overlapping boxes; boxes are regressed as
        relative offsets and log size ratios from a reference box.
  \item R-CNN classified warped proposals; Fast R-CNN shared the CNN and pooled RoI features;
        Faster R-CNN learned the proposals with anchors and an RPN.
  \item RoI Align replaces rounding by bilinear interpolation; Mask R-CNN adds a per-class $28\times28$
        sigmoid mask branch.
  \item One-stage detectors (YOLO, SSD, RetinaNet with the focal loss) predict densely; YOLACT and
        SOLO segment instances in one stage. Detectors are evaluated by AP over IoU thresholds.
\end{itemize}
\end{chaptersummary}

\section*{Exercises}
\addcontentsline{toc}{section}{Exercises}

\begin{exercise}\label{detection-segmentation:exr:1}
Two boxes of size $w\times w$ are shifted horizontally by $\delta w$, $0\leq\delta\leq1$. Show that their IoU is $(1 - \delta)/(1 + \delta)$.
For which shift is it $0.5$?
\solutionlink{detection-segmentation:exr:1}
\end{exercise}

\begin{exercise}\label{detection-segmentation:exr:2}
Run NMS on the boxes of \cref{detection-segmentation:ex:nms} with $\tau = 0.7$, and with the scores of $b_1$ and $b_2$
exchanged.
\solutionlink{detection-segmentation:exr:2}
\end{exercise}

\begin{exercise}\label{detection-segmentation:exr:3}
Decode the box predicted from the anchor $(x_a, y_a, w_a, h_a) = (100, 80, 64, 32)$ with
$t = (0.1, -0.25, \log2, 0)$. Why is $t_w = \log(w/w_a)$ a better regression target than $t_w = w - w_a$?
\solutionlink{detection-segmentation:exr:3}
\end{exercise}

\begin{exercise}\label{detection-segmentation:exr:4}
Show that bilinear interpolation reproduces exactly any function of the form $f(x,y) = a + bx + cy + dxy$, and
compute the derivative of the interpolated value with respect to $f_{11}$ and to $\alpha$.
\solutionlink{detection-segmentation:exr:4}
\end{exercise}
\end{notechapter}
